\chapter*{Abstract}
\addcontentsline{toc}{chapter}{Abstract}

This thesis examines monthly responses to high-frequency ECB monetary-policy surprises across financial, housing-credit, lending-condition, credit, and compensation-linked channels in Germany and the Euro Area. It asks whether post-COVID ECB easing shocks transmitted more strongly through housing-credit and selected financial-intermediation channels than through compensation-linked real-economy variables. The empirical design treats announcement-window surprises as policy-news innovations, aggregates them to the monthly level, and reports dynamic local-projection responses with explicit attention to identification sensitivity, response persistence, and uncertainty.

The findings indicate heterogeneous transmission. House-purchase lending growth is the clearest and most persistent response, while lending spreads and other intermediation variables provide supporting evidence for a housing-credit channel. Financial-market responses are mixed: the negative real DAX response is interpreted cautiously as evidence that ECB announcements may combine accommodation with adverse macroeconomic information. Compensation-linked proxies respond less consistently than the housing-credit block. The interpretation remains disciplined: the thesis provides reduced-form evidence on channel-specific responses while excluding household-level welfare, inequality, redistribution, and structural-mediation claims that the design does not identify.

\cleardoublepage
